%%%%%%%%%%%%%%%%%
%%% gen 11 %%%
%%%%%%%%%%%%%%%%

\Chapter{11}

\Verse{1}
{
	\hJ{וַיְהִי כל הָאָרֶץ שָׂפָה אֶחָת וּדְבָרִים אֲחָדִים׃}
}
{
	\eJ{
		And it was: all the earth was one language and the same
		words.
	}
}
{
	Here we go baby.

	After all those babies.
	
	The whole world is of one language, and one speech.
	
	An omnilingual orthography and auditory vocalization
	for the whole ass people of whole world as one nation.
	
	Finally something worth believing in.
}

\Verse{2}
{
	\hJ{וַיְהִי בְּנסְעָם מִקֶּדֶם וַיִּמְצְאוּ בִקְעָה בְּאֶרֶץ שִׁנְעָר וַיֵּשְׁבוּ שָׁם׃}
}
{
	\eJ{
		And it was when they were traveling from the east: and
		they found a valley in the land of Shinar and lived there.
	}
}
{
	Ok so they come from out east back to Sumer and stay there.
}

\Verse{3}
{
	\hJ{וַיֹּאמְרוּ אִישׁ אֶל רֵעֵהוּ הָבָה נִלְבְּנָה לְבֵנִים וְנִשְׂרְפָה לִשְׂרֵפָה וַתְּהִי לָהֶם הַלְּבֵנָה לְאָבֶן וְהַחֵמָר הָיָה לָהֶם לַחֹמֶר׃}
}
{
	\eJ{
		And they said to one another, “Come on, let’s make bricks
		and fire them.” And they had brick for stone, and they had
		bitumen for mortar.
	}
}
{
	And the One People say to each other,
	
	``Hey. We know how to build stuff.
}

\Verse{4}
{
	\hJ{וַיֹּאמְרוּ הָבָה נִבְנֶה לָּנוּ עִיר וּמִגְדָּל וְרֹאשׁוֹ בַשָּׁמַיִם}

	\hJ{וְנַעֲשֶׂה לָּנוּ שֵׁם פֶּן נָפוּץ עַל פְּנֵי כל הָאָרֶץ׃}
}
{
	\eJ{
		And they said, “Come on, let’s build ourselves a city and
		a tower that reaches to the heavens.
		
		And we’ll make ourselves a name, or else we’ll scatter over the face of
		all the earth.”
	}
}
{
	Let's build a big thing.
	
	Also let's decide on a name.
	
	I mean, I like you guys.

	We should make sure We don't forget each other.''

%	\footnote{
%		we’ll make ourselves a name. The story of the tower of Babylon has its
%		ambiguity at its center. It tells of an attempt by humanity to build a
%		tower whose top is in the sky, but it never tells what they will do when
%		they finish the tower, what is meant by “make ourselves a name,” or why
%		they fear being scattered if they do not have this tower. When YHWH sees
%		the tower He is described as concerned that “now nothing that they’ll
%		scheme to do will be precluded from them.” Their actions appear to be an
%		act of rebellion of some sort, and in this sense the episode fits with
%		those of Adam and Eve, Cain, and the flood, i.e., that humans as a
%		species are continuously in conflict with the initial “good” state of
%		creation. And so God disperses them and gives them languages that make
%		them less intelligible to one another and thereby less united. The story
%		thus prepares the way for a shift in the narrative away from dealing
%		with the fate of the species and, instead, dealing with individuals. At
%		the same time, this narrative provides the etiology of languages and of
%		the dispersion of humans all over the earth. It also provides a Hebrew
%		etiology for the name Babylon (Hebrew bbl) “because YHWH babbled (Hebrew
%		Ml) the language of all the earth” (11:9). This would have come as a
%		surprise to the people of Babylon, who understood the name to derive
%		from bab-ilu, the “Gate of God.”
%	}
}

\Verse{5}
{
	\hJ{וַיֵּרֶד יְהֹוָה לִרְאֹת אֶת הָעִיר וְאֶת הַמִּגְדָּל אֲשֶׁר בָּנוּ בְּנֵי הָאָדָם׃}
}
{
	\eJ{
		And YHWH went down to see the city and the tower that the
		children of humankind had built.
	}
}
{
	Then YHWH shows up.
	
	He looks at the One People.
	
	And the thing they're creating together.
}

\Verse{6}
{
	\hJ{וַיֹּאמֶר יְהֹוָה הֵן עַם אֶחָד וְשָׂפָה אַחַת לְכֻלָּם וְזֶה הַחִלָּם לַעֲשׂוֹת וְעַתָּה לֹא יִבָּצֵר מֵהֶם כֹּל אֲשֶׁר יָזְמוּ לַעֲשׂוֹת׃}
}
{
	\eJ{
		And YHWH said, “Here, they’re one people, and they all
		have one language, and this is what they’ve begun to do.
		And now nothing that they’ll scheme to do will be
		precluded from them.
	}
}
{
	And YHWH's like:
	
	Woah.\footnote{
		Lit. \heb{הֵן} in the original, an interjection expressing surprise.
	}
	
	The people are cooperating.
	
	Now they can just make plans and do things.
	
	Like anything.
	
	Together.
}

\Verse{7}
{
	\hJ{הָבָה נֵרְדָה וְנָבְלָה שָׁם שְׂפָתָם אֲשֶׁר לֹא יִשְׁמְעוּ אִישׁ שְׂפַת רֵעֵהוּ׃}
}
{
	\eJ{
		Come on, let’s go down and babble their language there so
		that one won’t understand another’s language.”
	}
}
{
	Guys,\footnote{Lit. \heb{הָבָה} in the original text, a word used
		before first-person plural future-tense verbs to invite someone
		to join in an action.
	}
	let's%
	\aC{\footnote{
	Yhwh unexpectedly speaks in the plural here, specifically
	the first-person plural (i.e., ``We tense'') just as in
	Genesis 1:26 (\textsl{Let us make humanity}) and 3:22
	(\textsl{The human has become like one of us}). The speaker here is
	Yhwh, singular, but the use of the plural form is widely understood
	to be an address to the divine council of heavenly beings, cf. the
	``Sons of the Gods'' in Genesis 6. See also Deut. 32:8--9 and
	Psalm 82 for two classic depictions of the divine council in the
	text. Some scholars argue further that the division of humanity into
	languages and nations in Gen. 11 is the same primordial event
	reflected in the Dead Sea Scrolls reading of Deuteronomy 32:8--9,
	where the Most High apportions the nations among the ``sons of God.''
	}}
	go mess everything up so they get confused and can't
	understand each other anymore.
}

\Verse{8}
{
	\hJ{וַיָּפֶץ יְהֹוָה אֹתָם מִשָּׁם עַל פְּנֵי כל הָאָרֶץ וַיַּחְדְּלוּ לִבְנֹת הָעִיר׃}
}
{
	\eJ{
		And YHWH scattered them from there over the face of all
		the earth, and they stopped building the city.
	}
}
{
	So YHWH invents nations and languages.

	For the people. Out of spite.
	
	And the people stop building the big thing.
}

\Verse{9}
{
	\hJ{עַל כֵּן קָרָא שְׁמָהּ בָּבֶל כִּי שָׁם בָּלַל יְהֹוָה שְׂפַת כּל הָאָרֶץ וּמִשָּׁם הֱפִיצָם יְהֹוָה עַל פְּנֵי כּל הָאָרֶץ׃}
}
{
	\eJ{
		On account of this its name was called Babylon, because
		YHWH babbled the language of all the earth there, and YHWH
		scattered them from there over the face of all the earth.
	}
}
{
	So it's called BBL, because BLL.

	The BBL, ladies and gentlemen.

	What?

	It's nothing.

	I just have something in my eye.
}

\Verse{10}
{
	\hR{אֵלֶּה תּוֹלְדֹת שֵׁם שֵׁם בֶּן מְאַת שָׁנָה וַיּוֹלֶד אֶת אַרְפַּכְשָׁד שְׁנָתַיִם אַחַר הַמַּבּוּל׃}
}
{
	\eR{
		These are the records of Shem: Shem was a hundred years
		old, and he fathered Arpachshad, two years after the
		flood.
	}
}
{
	I don't know what any of these fucking names mean.
	
	Bunch of gibberish from Giberia and babbling to me.
}

\Verse{11}
{
	\hR{וַיְחִי שֵׁם אַחֲרֵי הוֹלִידוֹ אֶת אַרְפַּכְשָׁד חֲמֵשׁ מֵאוֹת שָׁנָה וַיּוֹלֶד בָּנִים וּבָנוֹת׃}
}
{
	\eR{
		And Shem lived after his fathering Arpachshad five hundred
		years, and he fathered sons and daughters.
	}
}
{
	Arpachshad's not a name. That's not even, like, a name.
}

\Verse{12}
{
	\hR{וְאַרְפַּכְשַׁד חַי חָמֵשׁ וּשְׁלֹשִׁים שָׁנָה וַיּוֹלֶד אֶת שָׁלַח׃}
}
{
	\eR{
		And Arpachshad lived thirty-five years, and he fathered
		She-lah.
	}
}
{
	Screw these foreigners, they make no sense.
}

\Verse{13}
{
	\hR{וַיְחִי אַרְפַּכְשַׁד אַחֲרֵי הוֹלִידוֹ אֶת שֶׁלַח שָׁלֹשׁ שָׁנִים וְאַרְבַּע מֵאוֹת שָׁנָה וַיּוֹלֶד בָּנִים וּבָנוֹת׃}
}
{
	\eR{
		And Arpachshad lived after his fathering Shelah three
		years and four hundred years, and he fathered sons and
		daughters.
	}
}
{
	So blahbitty bloo lives a bunch of years after begatting gibberish.
}

\Verse{14}
{
	\hR{וְשֶׁלַח חַי שְׁלֹשִׁים שָׁנָה וַיּוֹלֶד אֶת עֵבֶר׃}
}
{
	\eR{And Shelah lived thirty years, and he fathered Eber.}
}
{
	Can we get all these foreigners out of the bible please?
}

\Verse{15}
{
	\hR{וַיְחִי שֶׁלַח אַחֲרֵי הוֹלִידוֹ אֶת עֵבֶר שָׁלֹשׁ שָׁנִים וְאַרְבַּע מֵאוֹת שָׁנָה וַיּוֹלֶד בָּנִים וּבָנוֹת׃}
}
{
	\eR{
		And Shelah lived after his fathering Eber three years and
		four hundred years, and he fathered sons and daughters.
	}
}
{
	Can. Not. Understand. You.%
	\footnote{Lit. 曰 in the original. Some manuscripts have Yue.}
}

\Verse{16}
{
	\hR{וַיְחִי עֵבֶר אַרְבַּע וּשְׁלֹשִׁים שָׁנָה וַיּוֹלֶד אֶת פָּלֶג׃}
}
{
	\eR{And Eber lived forty-three years, and he fathered Peleg.}
}
{
	Someone fathers peg-leg.
}

\Verse{17}
{
	\hR{וַיְחִי עֵבֶר אַחֲרֵי הוֹלִידוֹ אֶת פֶּלֶג שְׁלֹשִׁים שָׁנָה וְאַרְבַּע מֵאוֹת שָׁנָה וַיּוֹלֶד בָּנִים וּבָנוֹת׃}
}
{
	\eR{
		And Eber lived after his fathering Peleg thirty years and
		four hundred years, and he fathered sons and daughters.
	}
}
{
	Nobody lives this long, I don't care.
}

\Verse{18}
{
	\hR{וַיְחִי פֶלֶג שְׁלֹשִׁים שָׁנָה וַיּוֹלֶד אֶת רְעוּ׃}
}
{
	\eR{And Peleg lived thirty years, and he fathered Reu.}
}
{
	Why do we need to hear about these foreigners?
}

\Verse{19}
{
	\hR{וַיְחִי פֶלֶג אַחֲרֵי הוֹלִידוֹ אֶת רְעוּ תֵּשַׁע שָׁנִים וּמָאתַיִם שָׁנָה וַיּוֹלֶד בָּנִים וּבָנוֹת׃}
}
{
	\eR{
		And Peleg lived after his fathering Reu nine years and two
		hundred years, and he fathered sons and daughters.
	}
}
{
	Putting this on silent.
}

\Verse{20}
{
	\hR{וַיְחִי רְעוּ שְׁתַּיִם וּשְׁלֹשִׁים שָׁנָה וַיּוֹלֶד אֶת שְׂרוּג׃}
}
{
	\eR{And Reu lived thirty-two years, and he fathered Serug.}
}
{
	Serug isn't a name.
}

\Verse{21}
{
	\hR{וַיְחִי רְעוּ אַחֲרֵי הוֹלִידוֹ אֶת שְׂרוּג שֶׁבַע שָׁנִים וּמָאתַיִם שָׁנָה וַיּוֹלֶד בָּנִים וּבָנוֹת׃}
}
{
	\eR{
		And Reu lived after his fathering Serug seven years and
		two hundred years, and he fathered sons and daughters.
	}
}
{
	Ever feel like...
	
	We can't hear each other anymore.

	Sorry I'm going through a lot in this chapter. I just...

	I'll see you guys later.
}

\Verse{22}
{
	\hR{וַיְחִי שְׂרוּג שְׁלֹשִׁים שָׁנָה וַיּוֹלֶד אֶת נָחוֹר׃}
}
{
	\eR{And Serug lived thirty years, and he fathered Nahor.}
}
{
	\aC{
		Unlike many of the names in the preceding chapter, Nahor does
		not immediately suggest an obvious Hebrew etymology.
	}
}

\Verse{23}
{
	\hR{וַיְחִי שְׂרוּג אַחֲרֵי הוֹלִידוֹ אֶת נָחוֹר מָאתַיִם שָׁנָה וַיּוֹלֶד בָּנִים וּבָנוֹת׃}
}
{
	\eR{
		And Serug lived after his fathering Nahor two hundred
		years, and he fathered sons and daughters.
	}
}
{
	\aC{
		The genealogy continues without pausing to explain the
		names or their significance.
	}
}

\Verse{24}
{
	\hR{וַיְחִי נָחוֹר תֵּשַׁע וְעֶשְׂרִים שָׁנָה וַיּוֹלֶד אֶת תָּרַח׃}
}
{
	\eR{And Nahor lived twenty-nine years, and he fathered Terah.}
}
{
	\aC{
		Terah likewise lacks an evident Hebrew wordplay or
		transparent meaning.
	}
}

\Verse{25}
{
	\hR{וַיְחִי נָחוֹר אַחֲרֵי הוֹלִידוֹ אֶת תֶּרַח תְּשַׁע עֶשְׂרֵה שָׁנָה וּמְאַת שָׁנָה וַיּוֹלֶד בָּנִים וּבָנוֹת׃}
}
{
	\eR{
		And Nahor lived after his fathering Terah nineteen years
		and a hundred years, and he fathered sons and daughters.
	}
}
{
	\aC{
		The narrative's interest here is succession rather than
		etymology.
	}
}

\Verse{26}
{
	\hR{וַיְחִי תֶרַח שִׁבְעִים שָׁנָה וַיּוֹלֶד אֶת אַבְרָם אֶת נָחוֹר וְאֶת הָרָן׃}
}
{
	\eR{
		And Terah lived seventy years, and he fathered Abram,
		Nahor, and Haran.
	}
}
{
	\aC{
		The genealogy has now reached the family that will occupy
		the remainder of Genesis.
	}
}

\Verse{27}
{
	\hOther{וְאֵלֶּה תּוֹלְדֹת תֶּרַח תֶּרַח הוֹלִיד אֶת אַבְרָם אֶת נָחוֹר וְאֶת הָרָן וְהָרָן הוֹלִיד אֶת לוֹט׃}
}
{
	\eOther{
		And these are the records of Terah: Terah had fathered
		Abram, Nahor, and Haran, and Haran had fathered Lot.
	}
}
{
	\aC{
		The universal genealogy narrows here to a single household.
	}
}

\Verse{28}
{
	\hOther{וַיָּמת הָרָן עַל פְּנֵי תֶּרַח אָבִיו בְּאֶרֶץ מוֹלַדְתּוֹ בְּאוּר כַּשְׂדִּים׃}
}
{
	\eOther{
		And Haran died in the lifetime of Terah, his father, in
		the land of his birthplace, in Ur of the Chaldees.
	}
}
{
	Ok so everyone's in Ur.
	
	Haran dies.
}

\Verse{29}
{
	\hOther{וַיִּקַּח אַבְרָם וְנָחוֹר לָהֶם נָשִׁים שֵׁם אֵשֶׁת אַבְרָם שָׂרָי וְשֵׁם אֵשֶׁת נָחוֹר מִלְכָּה בַּת הָרָן אֲבִי מִלְכָּה וַאֲבִי יִסְכָּה׃}
}
{
	\eOther{
		And Abram and Nahor took wives. Abram's wife's name was
		Sarai, and Nahor's wife's name was Milcah, daughter of
		Haran, father of Milcah and father of Iscah.
	}
}
{
	People get married.
}

\Verse{30}
{
	\hP{וַתְּהִי שָׂרַי עֲקָרָה אֵין לָהּ וָלָד׃}
}
{
	\eP{And Sarai was infertile. She did not have a child.}
}
{
	Sarai can't have kids.
}

\Verse{31}
{
	\hP{וַיִּקַּח תֶּרַח אֶת אַבְרָם בְּנוֹ וְאֶת לוֹט בֶּן הָרָן בֶּן בְּנוֹ וְאֵת שָׂרַי כַּלָּתוֹ אֵשֶׁת אַבְרָם בְּנוֹ וַיֵּצְאוּ אִתָּם מֵאוּר כַּשְׂדִּים לָלֶכֶת אַרְצָה כְּנַעַן וַיָּבֹאוּ עַד חָרָן וַיֵּשְׁבוּ שָׁם׃}
}
{
	\eP{
		And Terah took Abram, his son, and Lot, son of Haran, his
		grandson, and Sarai, his daughter-in-law, the wife of
		Abram, his son; and they went with them from Ur of the
		Chaldees to go to the land of Canaan. And they came as far
		as Haran, and they stayed there.
	}
}
{
	Haran is dead, and they miss him, so they move to Haran.
	
	Like, out of nostalgia, or something.

	They leave Ur and go to Haran, and they stay there.
}

\Verse{32}
{
	\hR{וַיִּהְיוּ יְמֵי תֶרַח חָמֵשׁ שָׁנִים וּמָאתַיִם שָׁנָה וַיָּמת תֶּרַח בְּחָרָן׃}
}
{
	\eR{
		And Terah's days were five years and two hundred years.
		And Terah died in Haran.
	}
}
{
	Then Abe's dad dies.

	They're still in Haran.
	
	The town, not the corpse.
}